% Part: first-order-logic
% Chapter: completeness
% Section: complete-sets

% Definition of complete consistent sets. Properties of complete sets required
% for completeness proved are in provability.tex in the
% chapter on the proof system used.

\documentclass[../../../include/open-logic-section]{subfiles}

\begin{document}

\iftag{FOL}
      {\olfileid{fol}{com}{ccs}}
      {\olfileid{pl}{com}{ccs}}

\olsection{\usetoken{P}{sentence} चे संपूर्ण सुसंगत संच}

\begin{defn}[संपूर्ण संच]
\ollabel{def:complete-set} !!{sentence}s चा संच~$\Gamma$
\emph{!!{complete}} असतो तेव्हा आणि केवळ तेव्हाच, जेव्हा कोणत्याही
!!{sentence}~$!A$ साठी $!A \in \Gamma$ किंवा $\lnot !A \in \Gamma$.
\end{defn}

\begin{explain}
!!^{complete} वाक्यसंच कोणताही प्रश्न अनुत्तरित ठेवत नाहीत. कोणत्याही
!!{sentence}~$!A$ साठी $!A$ सत्य आहे की असत्य हे~$\Gamma$ “सांगतो.”
!!{complete} संचांचे महत्त्व संपूर्णता प्रमेयाच्या सिद्धतेपलीकडेही आहे.
उदाहरणार्थ, जी उपपत्ती !!{complete} असून स्वयंसिद्धकीय रीतीने मांडता
येते, तिच्याविषयी नेहमी निर्णय घेता येतो.
\end{explain}

\begin{explain}
संपूर्णतेच्या सिद्धतेत !!^{complete} सुसंगत संच महत्त्वाचे असतात, कारण
!!{sentence}s चा प्रत्येक सुसंगत संच~$\Gamma$ हा एखाद्या !!a{complete}
सुसंगत संच~$\Gamma^*$ मध्ये समाविष्ट असतो, याची हमी आपण देऊ शकतो.
!!^a{complete} सुसंगत संचात प्रत्येक !!{sentence}~$!A$ साठी $!A$ किंवा
त्याचे नकरण $\lnot !A$ यांपैकी एक असते, पण दोन्ही नसतात. हे विशेषतः
\iftag{FOL}{आण्विक !!{sentence}s}{विधानीय चले} यांच्यासाठी खरे असते.
म्हणून \iftag{FOL}{!!{constant}s ची योग्य भर घालून विस्तारित केलेल्या
भाषेतील}{} !!a{complete} सुसंगत संचापासून आपण अशी
\iftag{FOL}{!!a{structure}}{!!a{valuation}} रचू शकतो की तिच्यातील
\iftag{FOL}{!!{predicate}s चा अर्थ}{विधानीय चलांना नेमलेले सत्यतामूल्य}
हा~$\Gamma^*$ मध्ये कोणती \iftag{FOL}{आण्विक
!!{sentence}s}{!!{propositional variable}s} आहेत त्यानुसार परिभाषित
केला जातो. त्यानंतर ही \iftag{FOL}{!!{structure}}{!!{valuation}}
$\Gamma^*$ मधील सर्व !!{sentence}s (आणि म्हणून~$\Gamma$ मधीलही सर्व)
सत्य करते, हे दाखवता येते. या शेवटच्या तथ्याच्या सिद्धतेसाठी
$\lnot !A \in \Gamma^*$ तेव्हा आणि केवळ तेव्हाच, जेव्हा $!A \notin
\Gamma^*$; $(!A \lor !B) \in \Gamma^*$ तेव्हा आणि केवळ तेव्हाच, जेव्हा
$!A \in \Gamma^*$ किंवा $!B \in \Gamma^*$; इत्यादी तथ्ये आवश्यक असतात.
\end{explain}

पुढे आपण $\Proves$ ची परावर्तिता, एकस्वनिकता आणि संक्रमकता हे गुणधर्म
अनेकदा अध्याहृतपणे वापरू (पाहा
\iftag{FOL}{%
  \tagrefs{prfSC/{fol:seq:ptn:sec},prfND/{fol:ntd:ptn:sec},prfAX/{fol:axd:ptn:sec},prfTab/{fol:tab:ptn:sec}}}{%
  \tagrefs{prfSC/{pl:seq:ptn:sec},prfND/{pl:ntd:ptn:sec},prfAX/{pl:axd:ptn:sec},prfTab/{pl:tab:ptn:sec}}}).

\begin{prop}
\ollabel{prop:ccs}
$\Gamma$ हा !!{complete} आणि सुसंगत आहे असे समजा. मग:
\begin{enumerate}
\item \ollabel{prop:ccs-prov-in} जर $\Gamma \Proves !A$, तर $!A \in
  \Gamma$.

\tagitem{prvAnd}{\ollabel{prop:ccs-and} $!A \land !B \in \Gamma$
  तेव्हा आणि केवळ तेव्हाच, जेव्हा $!A \in \Gamma$ आणि $!B \in \Gamma$
  ही दोन्ही खरे आहेत.}{}

\tagitem{prvOr}{\ollabel{prop:ccs-or} $!A \lor !B \in \Gamma$ तेव्हा
  आणि केवळ तेव्हाच, जेव्हा $!A \in \Gamma$ किंवा $!B \in \Gamma$.}{}

\tagitem{prvIf}{\ollabel{prop:ccs-if} $!A \lif !B \in \Gamma$ तेव्हा
  आणि केवळ तेव्हाच, जेव्हा $!A \notin \Gamma$ किंवा $!B \in \Gamma$.}{}
\end{enumerate}
\end{prop}

\begin{proof}
पुढील सर्व प्रकरणांत $\Gamma$ हा !!{complete} आणि सुसंगत आहे असे समजू.
\begin{enumerate}
\item जर $\Gamma \Proves !A$, तर $!A \in \Gamma$.

$\Gamma \Proves !A$ असे समजा. उलटपक्षी $!A \notin \Gamma$ असे गृहीत
धरा. $\Gamma$ हा !!{complete} असल्यामुळे $\lnot !A \in \Gamma$. पुढील
संदर्भांनुसार
\iftag{FOL}{%
  \tagrefs{prfSC/{fol:seq:prv:prop:explicit-inc},
    prfND/{fol:ntd:prv:prop:explicit-inc},
    prfAX/{fol:axd:prv:prop:explicit-inc},
    prfTab/{fol:tab:prv:prop:explicit-inc}}}{%
  \tagrefs{prfSC/{pl:seq:prv:prop:explicit-inc},
    prfND/{pl:ntd:prv:prop:explicit-inc},
    prfAX/{pl:axd:prv:prop:explicit-inc},
    prfTab/{pl:tab:prv:prop:explicit-inc}}},
$\Gamma$ विसंगत आहे. हे $\Gamma$ सुसंगत असल्याच्या गृहीतकाशी विसंगत आहे.
म्हणून $!A \notin \Gamma$ असणे शक्य नाही; त्यामुळे $!A \in \Gamma$.

\tagitem{defAnd}{}{%
\iftag{probAnd}{सराव.}{%
$!A \land !B \in \Gamma$ तेव्हा आणि केवळ तेव्हाच, जेव्हा $!A \in
\Gamma$ आणि $!B \in \Gamma$ ही दोन्ही खरे आहेत:

पुढे जाणाऱ्या दिशेसाठी $!A \land !B \in \Gamma$ असे समजा. मग
पुढील संदर्भांतील
\iftag{FOL}{%
  \tagrefs{prfSC/{fol:seq:ppr:prop:provability-land},
    prfND/{fol:ntd:ppr:prop:provability-land},
    prfAX/{fol:axd:ppr:prop:provability-land},
    prfTab/{fol:tab:ppr:prop:provability-land}}}{%
  \tagrefs{prfSC/{pl:seq:ppr:prop:provability-land},
    prfND/{pl:ntd:ppr:prop:provability-land},
    prfAX/{pl:axd:ppr:prop:provability-land},
    prfTab/{pl:tab:ppr:prop:provability-land}}}, बाब~(1) नुसार
$\Gamma \Proves !A$ आणि $\Gamma \Proves !B$. \olref{prop:ccs-prov-in}
नुसार $!A \in \Gamma$ आणि $!B \in \Gamma$, अपेक्षेप्रमाणे.

उलट दिशेसाठी $!A \in \Gamma$ आणि $!B \in \Gamma$ असे समजा. पुढील
संदर्भांतील
\iftag{FOL}{%
  \tagrefs{prfSC/{fol:seq:ppr:prop:provability-land},
    prfND/{fol:ntd:ppr:prop:provability-land},
    prfAX/{fol:axd:ppr:prop:provability-land},
    prfTab/{fol:tab:ppr:prop:provability-land}}}{%
  \tagrefs{prfSC/{pl:seq:ppr:prop:provability-land},
    prfND/{pl:ntd:ppr:prop:provability-land},
    prfAX/{pl:axd:ppr:prop:provability-land},
    prfTab/{pl:tab:ppr:prop:provability-land}}}, बाब~(2) नुसार
$\Gamma \Proves !A \land !B$. \olref{prop:ccs-prov-in} नुसार $!A \land
!B \in \Gamma$.}}

\tagitem{defOr}{}{%
\iftag{probOr}{सराव.}{%
प्रथम $!A \lor !B \in \Gamma$ असेल, तर $!A \in \Gamma$ किंवा $!B \in
\Gamma$, हे दाखवू. $!A \lor !B \in \Gamma$, पण $!A \notin \Gamma$ आणि
$!B \notin \Gamma$ असे समजा. $\Gamma$ हा !!{complete} असल्यामुळे
$\lnot !A \in \Gamma$ आणि $\lnot !B \in \Gamma$. पुढील संदर्भांतील
\iftag{FOL}{%
  \tagrefs{prfSC/{fol:seq:ppr:prop:provability-lor},
    prfND/{fol:ntd:ppr:prop:provability-lor},
    prfAX/{fol:axd:ppr:prop:provability-lor},
    prfTab/{fol:tab:ppr:prop:provability-lor}}}{%
  \tagrefs{prfSC/{pl:seq:ppr:prop:provability-lor},
    prfND/{pl:ntd:ppr:prop:provability-lor},
    prfAX/{pl:axd:ppr:prop:provability-lor},
    prfTab/{pl:tab:ppr:prop:provability-lor}}},
बाब~(1) नुसार $\Gamma$ विसंगत आहे; हा विरोधाभास आहे. म्हणून $!A \in
\Gamma$ किंवा $!B \in \Gamma$.

उलट दिशेसाठी $!A \in \Gamma$ किंवा $!B \in \Gamma$ असे समजा. पुढील
संदर्भांतील
\iftag{FOL}{%
  \tagrefs{prfSC/{fol:seq:ppr:prop:provability-lor},
    prfND/{fol:ntd:ppr:prop:provability-lor},
    prfAX/{fol:axd:ppr:prop:provability-lor},
    prfTab/{fol:tab:ppr:prop:provability-lor}}}{%
  \tagrefs{prfSC/{pl:seq:ppr:prop:provability-lor},
    prfND/{pl:ntd:ppr:prop:provability-lor},
    prfAX/{pl:axd:ppr:prop:provability-lor},
    prfTab/{pl:tab:ppr:prop:provability-lor}}}, बाब~(2) नुसार
$\Gamma \Proves !A \lor !B$. \olref{prop:ccs-prov-in} नुसार $!A \lor
!B \in \Gamma$, अपेक्षेप्रमाणे.}}

\tagitem{defIf}{}{%
\iftag{probIf}{सराव.}{%
पुढे जाणाऱ्या दिशेसाठी $!A \lif !B \in \Gamma$ असे समजा आणि उलटपक्षी
$!A \in \Gamma$ व $!B \notin \Gamma$ असे गृहीत धरा. या गृहीतकांनुसार
$!A \lif !B \in \Gamma$ आणि $!A \in \Gamma$. पुढील संदर्भांतील
\iftag{FOL}{%
  \tagrefs{prfSC/{fol:seq:ppr:prop:provability-lif},
    prfND/{fol:ntd:ppr:prop:provability-lif},
    prfAX/{fol:axd:ppr:prop:provability-lif},
    prfTab/{fol:tab:ppr:prop:provability-lif}}}{%
  \tagrefs{prfSC/{pl:seq:ppr:prop:provability-lif},
    prfND/{pl:ntd:ppr:prop:provability-lif},
    prfAX/{pl:axd:ppr:prop:provability-lif},
    prfTab/{pl:tab:ppr:prop:provability-lif}}}, बाब~(1) नुसार
$\Gamma \Proves !B$. पण मग \olref{prop:ccs-prov-in} नुसार $!B \in
\Gamma$; हे $!B \notin \Gamma$ या गृहीतकाशी विसंगत आहे.

उलट दिशेसाठी प्रथम $!A \notin \Gamma$ हे प्रकरण पाहू. $\Gamma$ हा
!!{complete} असल्यामुळे $\lnot !A \in \Gamma$. पुढील संदर्भांतील
\iftag{FOL}{%
  \tagrefs{prfSC/{fol:seq:ppr:prop:provability-lif},
    prfND/{fol:ntd:ppr:prop:provability-lif},
    prfAX/{fol:axd:ppr:prop:provability-lif},
    prfTab/{fol:tab:ppr:prop:provability-lif}}}{%
  \tagrefs{prfSC/{pl:seq:ppr:prop:provability-lif},
    prfND/{pl:ntd:ppr:prop:provability-lif},
    prfAX/{pl:axd:ppr:prop:provability-lif},
    prfTab/{pl:tab:ppr:prop:provability-lif}}}, बाब~(2) नुसार
$\Gamma \Proves !A \lif !B$. पुन्हा \olref{prop:ccs-prov-in} नुसार
$!A \lif !B \in \Gamma$, अपेक्षेप्रमाणे.

आता $!B \in \Gamma$ हे प्रकरण पाहू. पुढील संदर्भांतील
\iftag{FOL}{%
  \tagrefs{prfSC/{fol:seq:ppr:prop:provability-lif},
    prfND/{fol:ntd:ppr:prop:provability-lif},
    prfTab/{fol:tab:ppr:prop:provability-lif},
    prfAX/{fol:axd:ppr:prop:provability-lif}}}{%
  \tagrefs{prfSC/{pl:seq:ppr:prop:provability-lif},
    prfND/{pl:ntd:ppr:prop:provability-lif},
    prfTab/{pl:tab:ppr:prop:provability-lif},
    prfAX/{pl:axd:ppr:prop:provability-lif}}},
पुन्हा बाब~(2) नुसार $\Gamma \Proves !A \lif !B$.
\olref{prop:ccs-prov-in} नुसार $!A \lif !B \in \Gamma$.}}
\end{enumerate}
\end{proof}

\tagprob{FOL}
\begin{prob}
\olref[fol][com][ccs]{prop:ccs} ची सिद्धता पूर्ण करा.
\end{prob}
\tagendprob

\tagprob{notFOL}
\begin{prob}
\olref[pl][com][ccs]{prop:ccs} ची सिद्धता पूर्ण करा.
\end{prob}
\tagendprob

\end{document}
